Otros dos conceptos importantes a la hora de trabajar con algoritmos de inferencia de clases, son el overfitting y el underfitting. El overfitting ocurre cuando un modelo se ajusta demasiado a los datos de entrenamiento, capturando ruido o patrones específicos que no generalizan bien a nuevos datos. Esto puede resultar en un rendimiento excelente en el conjunto de entrenamiento pero pobre en datos no vistos.
Por otro lado, el underfitting sucede para el caso contrario.cuando un modelo es demasiado simple para capturar la complejidad de los datos, lo que lleva a un rendimiento deficiente tanto en el conjunto de entrenamiento como en datos nuevos.
Finalmente otro concepto importante es la generalización.Esta se refiere a la capacidad de un modelo para desempeñarse bien en datos no vistos, y es crucial encontrar un equilibrio entre overfitting y underfitting para lograr una buena generalización.